历史上，32-bit 的 $n_t d$ 在 $2^{31}$ 处溢出，导致主码尺度归零；后续进程又因 cache key 不含代码版本而加载了错误模型。分块计算范数并以 double 累积修复了溢出。这一故障要求分别核查模型几何、编码行与放置关系，即使粗路由成功也不例外。

\subsection{已完成的表示证据}

10 月 13 日的残差下界 pilot 从 arms6 训练缓存导出 $\Pi$、码本、ID、偏移、修正项及逻辑解码后的打包编码，覆盖完整的 Vogue（932,328 行）、arXiv（2,253,000 行）和 OpenAI-3072（999,000 行）数据库。

码本检查发现可分离性不匹配、标量级值不匹配、残差码本无序对均为零。独立解码导出的主码／残差后，FP64 循环遍历每一行，核对 $\gamma=2\langle\widehat p,\widehat r\rangle$。三个数据库的最大修正项绝对误差为$1.28\times10^{-7}$--$2.77\times10^{-7}$；采样GPU返回top-10分数的最大相对误差为$2.38\times10^{-7}$--$4.85\times10^{-5}$，定义为$|s_{\rm GPU}-s_{64}|/s_{64}$，参考含查询范数常数。这些是导出表示的检查，不是从原始输入向量重新编码；逐数据集诊断保留在仓库补充材料。


在七个工作负载、896 个查询实例中，pilot 对 72,800,590 个路由候选实例检查所提残差下界与 FP64 重建距离，未报告违反；同一集合也支持 FP64 组合分数恒等式检查。GPU 返回分数检查仅覆盖 top-10 槽位，并跳过路由映射中不存在的 ID，因此尚未完整验证输出成员关系。

规范化命名空间并去掉导出方法后，编码、主码／残差训练与构建源码同评测实现一致（\path{scripts/audit_ch4_review.py}）。该对应覆盖共同路径；后续模型身份与大数据集两遍构建仍需单独核查。

\subsection{独立输入与放置审计}
\label{sec:4-representation-audit}

补证导出器在Vogue、arXiv、OpenAI-3072与Stella上使用共同构建路径。
四者全行偏移均单调、端点为0和$N$，ID均构成数据库行的置换。
对存储旋转矩阵做FP64乘法，$\|\Pi\Pi^{\mathsf T}-I\|$最大逐项误差为
$3.96$--$4.34\times10^{-7}$。采样范数检查使用FP64旋转而非GPU GEMM，
因此验证的是存储矩阵的几何性质。\Cref{tab:4-representation-audit}汇总采样重编码。

\par\medskip\noindent\begin{minipage}{\linewidth}
\centering\captionsetup{hypcap=false}
\captionof{table}[原始输入 FP64 重编码审计]{原始输入FP64重编码审计。差异数按坐标决策计数，最后一列为残差不一致
坐标离决策边界的最大距离。结构检查覆盖全行，重编码采用采样。}
\label{tab:4-representation-audit}
\scriptsize\setlength{\tabcolsep}{3pt}
\begin{tabular}{@{}lrrrrr@{}}
\toprule
数据集 & 重编码行数 & 正交误差 & 主码差异 & 残差差异 & 最大边界距离 \\
\midrule
Vogue & 200,000 & 4.11e-07 & 14 & 5,409 & 1.12e-07 \\
arXiv & 100,000 & 4.11e-07 & 1 & 2,515 & 7.31e-08 \\
OpenAI-3072 & 20,000 & 3.96e-07 & 5,745 & 1,646,806 & 2.67e-05 \\
Stella & 100,000 & 4.34e-07 & 16 & 3,064 & 1.03e-07 \\
\bottomrule
\end{tabular}

\end{minipage}\par\medskip

768／1,024维上的差异均在决策边界$1.12\times10^{-7}$内。
OpenAI-3072则明显不同：残差不一致占所查坐标的2.68\%，最远离边界
$2.67\times10^{-5}$，与TF32旋转一致，因此不能将所有数据集称为FP32等价。
残差比较以已存主码位为条件，避免重复计入主码翻转。

三个较小数据库在同一模型下强制一遍／两遍构建，按ID关联的主码和残差字节完全
相同，correction差至多$8.94\times10^{-8}$；列表分配分别有1、1、40行不同。
Stella只验证自然两遍路径。导出器不含IVF质心，因此列表成员决策本身仍未独立验证，
结构正确不等于已证明分配决策正确。

\subsection{契约覆盖与剩余检查}
% Keep this float after the preceding fixed-position recoding-audit table.
\suppressfloats[t]

\begin{table}[tbp]
\centering
\caption[表示契约与审计覆盖]{纳入pilot及补证输入审计后的表示契约，区分存储几何、采样算术与行放置的覆盖。}
\label{tab:4-contracts}
\footnotesize\setlength{\tabcolsep}{4pt}
\begin{tabularx}{\linewidth}{@{}>{\raggedright\arraybackslash}p{1.8cm}*{2}{>{\raggedright\arraybackslash}X}@{}}
\toprule
边界 & 所需性质 & 证据／剩余检查 \\
\midrule
C0：几何 & 共同变换，浮点等距误差有界 & 四数据集已查存储矩阵正交性；区分TF32 GPU算术 \\
C1：码本／修正 & 笛卡尔级值、有序残差码本，修正项匹配解码编码 & 三个数据库 pilot 已完成；全行修正误差见\cref{sec:4-6} \\
C1：输入编码 & 存储编码对应预期原始输入行 & 四份原始输入采样已做FP64重编码；量化边界与TF32差异已量化 \\
C1：放置 & 偏移覆盖行，ID 为置换，字段遵循同一置换 & 全行偏移／ID核对及三份按ID关联的路径比较；质心成员决策未验证 \\
执行 & 索引、同步和分数顺序契约有效 & 已有查询与 sanitizer 检查，有限覆盖见\cref{sec:7-4} \\
\bottomrule
\end{tabularx}
\end{table}

构建要求 $M\mid d$、$D_s\mid B$、$M\bmod4=0$，LUT 验证使用实际码字切片。
输入须有限、计数可表示并预留哨兵位置。ID／位置用 signed 32-bit；较大计数与地址乘积用更宽算术。

审计使用导出模型而非按 seed 重生成。T4 检查计算分数选择，与输入重编码分开
（\cref{sec:7-strict-selection}）；质心成员关系、空列表与对抗尾批次仍待验证。
